\chapter{Generazione del Java Bytecode}\label{chap:java-bytecode}
%
\vspace*{-2ex}
\begin{center}
\includegraphics[width=6.5cm]{cat9.jpg}
\end{center}
%\vspace*{-2ex}

Il Kitten bytecode \e molto semplice, poich\'e non \e necessario
preoccuparsi dell'esistenza di forme ottimizzate della stessa istruzione,
o dell'assenza di bytecode per alcune operazioni su alcuni tipi.
\E quindi ideale per la generazione del codice intermedio e la
sua eventuale ottimizzazione. A partire dal bytecode Kitten, \e poi possibile
generare del codice eseguibile. Questo codice eseguibile
potrebbe essere specifico di una data architettura su cui si intende
eseguire il programma, come ad esempio il linguaggio macchina
x86, oppure essere indipendente dalla macchina, basato quindi su
un bytecode astratto eseguibile su una qualche macchina virtuale.
In questo capitolo seguiremo la seconda strada ed esamineremo
la generazione del codice eseguibile in Java bytecode.

Il Java bytecode somiglia al Kitten bytecode, nel senso che le sue
istruzioni operano su uno stack degli operandi, su delle variabili
locali e su uno stack di attivazione. Anche le singole istruzioni e
la tecnica di chiamata e ritorno da metodo sono molto simili a quelle
del Kitten bytecode. \E per\`o un linguaggio meno esplicitamente
tipato, nel senso che, per ridurre l'occupazione in memoria,
molte istruzioni non riportano esplicitamente i
tipi su cui operano, quando questi tipi possono essere inferiti a partire
dal programma.
Il Java bytecode \e specificato in forma binaria all'interno di file
Java eseguibili detti \emph{file classe}.

La generazione del Java bytecode all'interno dei file classe \e tecnicamente
complessa per vari motivi:

\begin{enumerate}
\item a una singola istruzione del Kitten bytecode potrebbero corrispondere
      \piu istruzioni alternative del Java bytecode, variamente ottimizzate.
      Idealmente, occorre scegliere l'alternativa migliore in termini di
      occupazione in byte e di tempo di esecuzione;
\item ad alcune istruzioni del Kitten bytecode potrebbe non corrispondere
      una singola istruzione del Java bytecode, ma una sequenza di \piu
      istruzioni;
\item la struttura a blocchi del Kitten bytecode deve essere linearizzata in
      una sequenza di istruzioni, introducendo numeri di linea e,
      ove necessario, delle istruzioni di salto esplicito;
\item la sequenza di istruzioni deve infine essere trasformata in codice
      binario, facendo corrispondere a ogni istruzioni un byte, seguito da
      una qualche codifica dei suoi operandi, sempre in termini di byte;
\item i salti all'interno del Java bytecode devono venire specificati come offset
      di numeri di linea dal punto di partenza a quello di arrivo. Per
      salti particolarmente lunghi, per cui l'offset va in overflow,
      occorre creare dei \emph{punti di appoggio} intermedi;
\item le costanti del linguaggio (numeriche o stringhe) devono essere
      inserite all'interno di una tavola delle costanti del file classe.
      Nel caso in cui occorressero \piu volte nel codice, conviene
      riciclare la stessa costante piuttosto che inserirne due identiche
      nella tavola delle costanti.
\end{enumerate}
%
Tutti questi problemi giustificano l'utilizzo di una libreria di
supporto per la creazione e manipolazione del Java bytecode. Abbiamo
quindi scelto di utilizzare a tale scopo
la libreria BCEL (ByteCode Engineering Library).

\section{La generazione del Java bytecode per i bytecode Kitten sequenziali}

Abbiamo detto che a una singola istruzione del Kitten bytecode potrebbero
corrisponderne \piu di una del Java bytecode. Per questo motivo, il metodo
di traduzione da Kitten bytecode in Java bytecode per i bytecode sequenziali
\e definito dentro \texttt{bytecode/NonBranchingBytecode.java} come
%
\begin{verbatim}
  public abstract InstructionList generateJavaBytecode
                                   (JavaClassGenerator classGen);
\end{verbatim}
%
Il tipo di ritorno di tale metodo
\e \texttt{InstructionList}, \cioe una classe di BCEL che rappresenta
una sequenza (eventualmente vuota) di bytecode Java.
L'oggetto passato come argomento \e il generatore di Java bytecode, che vedremo
alla fine del capitolo. Per adesso, ci interessa solo sapere che al
suo interno esiste una \emph{fattoria di istruzioni}, \cioe una classe
di BCEL che aiuta il programmatore a generare Java bytecode. Ad esempio,
tale fattoria ci aiuter\`a nella scelta dell'istruzione \piu ottimizzata
quando saremo davanti a \piu alternative e inserir\`a automaticamente
delle costanti nella tavola delle costanti, per quei bytecode Java che
utilizzano delle costanti. Tale fattoria pu\`o essere ottenuta invocando
\texttt{classGen.getFactory()}.

\begin{figure}
\begin{center}
\begin{tabular}{l|l}
\multicolumn{1}{c}{Kitten} &
\multicolumn{1}{c}{Java} \\\hline\hline
$\mathtt{nop}$ & $\mathtt{nop}$ \\\hline
$\mathtt{pop}$ $t$ & $\mathtt{pop}$ \\\hline
$\mathtt{dup}$ $t$ & $\mathtt{dup}$ \\\hline
$\mathtt{const}\ \mathit{nil}$ & $\mathtt{aconst\_null}$\\\hline
$\mathtt{const}\ \mathit{float\_number}$ & $\mathtt{fconst}\ \mathit{float\_number}$\\\hline
$\mathtt{const -\!\!1,0,1,2,3}$ & $\mathtt{iconst -\!\!1,0,1,2,3}$\\\hline
$\mathtt{const}\ \mathit{8\_bits\_integer}$ & $\mathtt{bipush}\ \mathit{8\_bits\_integer}$\\\hline
$\mathtt{const}\ \mathit{16\_bits\_integer}$ & $\mathtt{sipush}\ \mathit{16\_bits\_integer}$\\\hline
$\mathtt{const}\ \mathit{32\_bits\_integer}$ & $\mathtt{ldc}\ \mathit{32\_bits\_integer}$\\\hline
$\mathtt{const}\ \mathit{true}$ & $\mathtt{iconst\ 1}$\\\hline
$\mathtt{const}\ \mathit{false}$ & $\mathtt{iconst\ 0}$\\\hline
$\mathtt{load\ \mathit{l}\ of\ type\ int,boolean}$ & $\mathtt{iload\ \mathit{l}}$\\\hline
$\mathtt{load\ \mathit{l}\ of\ type\ float}$ & $\mathtt{fload\ \mathit{l}}$\\\hline
$\mathtt{load\ \mathit{l}\ of\ type\ nil,\mathit{reference}}$ & $\mathtt{aload\ \mathit{l}}$\\\hline
$\mathtt{store\ \mathit{l}\ of\ type\ int,boolean}$ & $\mathtt{istore\ \mathit{l}}$\\\hline
$\mathtt{store\ \mathit{l}\ of\ type\ float}$ & $\mathtt{fstore\ \mathit{l}}$\\\hline
$\mathtt{store\ \mathit{l}\ of\ type\ nil,\mathit{reference}}$ & $\mathtt{astore\ \mathit{l}}$\\\hline
$\mathtt{neg\ int}$ & $\mathtt{ineg}$\\\hline
$\mathtt{neg\ float}$ & $\mathtt{fneg}$\\\hline
$\mathtt{neg\ boolean}$ & $\begin{array}{rl}
                            & \mathtt{ifeq}\ \mathit{after} \\
                            & \mathtt{iconst\ 0} \\
                            & \mathtt{goto}\ \mathit{end}\\
                            \mathit{after:} & \mathtt{iconst\ 1}\\
                            \mathit{end:} & \mathtt{nop}
                           \end{array}$\\\hline
$\mathtt{add\ int}$ & $\mathtt{iadd}$\\\hline
$\mathtt{add\ float}$ & $\mathtt{fadd}$\\\hline
$\mathtt{add/sub/div/mul\ int}$ & $\mathtt{iadd/isub/idiv/imul}$\\\hline
$\mathtt{add/sub/div/mul\ float}$ & $\mathtt{fadd/fsub/fdiv/fmul}$\\\hline
\end{tabular}
\end{center}
\caption{La traduzione da bytecode Kitten a bytecode Java (1/3).}
  \label{fig:kitten_into_java_1}
\end{figure}

\begin{figure}
\begin{center}
\begin{tabular}{l|l}
\multicolumn{1}{c}{Kitten} &
\multicolumn{1}{c}{Java} \\\hline\hline
$\mathtt{eq\ int}$ & $\begin{array}{rl}
                            & \mathtt{if\_icmpeq}\ \mathit{after} \\
                            & \mathtt{iconst\ 0} \\
                            & \mathtt{goto}\ \mathit{end}\\
                            \mathit{after:} & \mathtt{iconst\ 1}\\
                            \mathit{end:} & \mathtt{nop}
                           \end{array}$\\\hline
$\mathtt{eq\ float}$ & $\begin{array}{rl}
                            & \mathtt{fcmpl} \\
                            & \mathtt{ifeq}\ \mathit{after} \\
                            & \mathtt{iconst\ 0} \\
                            & \mathtt{goto}\ \mathit{end}\\
                            \mathit{after:} & \mathtt{iconst\ 1}\\
                            \mathit{end:} & \mathtt{nop}
                           \end{array}$\\\hline
$\mathtt{eq\ \mathit{reference}}$ & $\begin{array}{rl}
                            & \mathtt{if\_acmpeq}\ \mathit{after} \\
                            & \mathtt{iconst\ 0} \\
                            & \mathtt{goto}\ \mathit{end}\\
                            \mathit{after:} & \mathtt{iconst\ 1}\\
                            \mathit{end:} & \mathtt{nop}
                           \end{array}$\\\hline
$\mathtt{new\ }\mathit{\kappa}$ & $\mathtt{new\ }\mathit{\kappa}$ \\\hline
$\mathtt{newstring}\ \mathit{string}$ & $\begin{array}{rl}
                            & \mathtt{new\ runTime.String} \\
                            & \mathtt{dup} \\
                            & \mathtt{ldc}\ \mathit{string}\\
                            & \mathtt{invokespecial\ runTime.String.\langle init\rangle():void}
                           \end{array}$\\\hline
$\mathtt{newarray\ of\ \mathtt{int/float/boolean}}$ & $\mathtt{newarray\ \mathtt{int/float/boolean}}$ \\\hline
$\mathtt{newarray\ of\ \mathit{reference}}$ & $\mathtt{anewarray\ of\ \mathit{reference}}$ \\\hline
$\mathtt{cast\ int\ into\ float}$ & $\mathtt{i2f}$ \\\hline
$\mathtt{cast\ float\ into\ int}$ & $\mathtt{f2i}$ \\\hline
$\mathtt{cast}\ \mathit{reference}_1\ \mathtt{into}\ \mathit{reference}_2$ & $\mathtt{checkcast}\ \mathit{reference}_2$\\\hline
$\mathtt{getfield\ }\kappa.f$ & $\mathtt{getfield\ }\kappa.f$ \\\hline
$\mathtt{putfield\ }\kappa.f$ & $\mathtt{putfield\ }\kappa.f$ \\\hline
$\mathtt{arrayload\ from\ array\ of\ int}$ & $\mathtt{iaload}$ \\\hline
$\mathtt{arrayload\ from\ array\ of\ float}$ & $\mathtt{faload}$ \\\hline
$\mathtt{arrayload\ from\ array\ of\ boolean}$ & $\mathtt{baload}$ \\\hline
$\mathtt{arrayload\ from\ array\ of\ }\mathit{reference}$ & $\mathtt{aaload}$ \\\hline
$\mathtt{arraystore\ into\ array\ of\ int}$ & $\mathtt{iastore}$ \\\hline
$\mathtt{arraystore\ into\ array\ of\ float}$ & $\mathtt{fastore}$ \\\hline
$\mathtt{arraystore\ into\ array\ of\ boolean}$ & $\mathtt{bastore}$ \\\hline
$\mathtt{arraystore\ into\ array\ of\ }\mathit{reference}$ & $\mathtt{aastore}$ \\\hline
\end{tabular}
\end{center}
\caption{La traduzione da bytecode Kitten a bytecode Java (2/3).}
  \label{fig:kitten_into_java_2}
\end{figure}

\begin{figure}
\begin{center}
\begin{tabular}{l|l}
\multicolumn{1}{c}{Kitten} &
\multicolumn{1}{c}{Java} \\\hline\hline
$\mathtt{constructorcall\ \kappa(}$$\vec{t}$$\mathtt{):void}$ & $\mathtt{invokespecial\ \kappa.\langle init\rangle(}$$\vec{t}$$\mathtt{):void}$ \\\hline
$\mathtt{virtualcall\ \kappa.\mathit{m}(}$$\vec{t}$$\mathtt{):\mathit{t'}}$ & $\mathtt{invokevirtual\ \kappa.\mathit{m}(}$$\vec{t}$$\mathtt{):\mathit{t'}}$ \\\hline
$\mathtt{return\ int}$ & $\mathtt{ireturn}$ \\\hline
$\mathtt{return\ float}$ & $\mathtt{freturn}$ \\\hline
$\mathtt{return\ boolean}$ & $\mathtt{ireturn}$ \\\hline
$\mathtt{return\ void}$ & $\mathtt{return}$ \\\hline
\end{tabular}
\end{center}
\caption{La traduzione da bytecode Kitten a bytecode Java (3/3).}
  \label{fig:kitten_into_java_3}
\end{figure}

Esaminiamo alcuni esempi di ridefinizione di tale metodo all'interno delle
sottoclassi dei bytecode Kitten. Le Figure~\ref{fig:kitten_into_java_1},
\ref{fig:kitten_into_java_2} and~\ref{fig:kitten_into_java_3} riportano
la lista dei bytecode Kitten e la loro traduzione in bytecode Java. Il primo
esempio \e quello del bytecode Kitten $\mathtt{nop}$, a cui corrisponde un
identico Java bytecode. L'implementazione in \texttt{bytecode/NOP.java} \e quindi semplicissima:
%
\begin{verbatim}
  public InstructionList generateJavaBytecode
      (JavaClassGenerator classGen) {
    return new InstructionList(InstructionFactory.NOP);
  }
\end{verbatim}
%
ritorniamo \cioe una lista di una singola istruzione, una $\mathtt{nop}$. Usiamo qui un
campo statico
della fattoria BCEL che ci permette di fare riferimento a un'istruzione gi\`a creata, al posto
di costruirla ogni volta tramite un costruttore.

Nel caso del bytecode Kitten $\mathtt{pop}\ \mathit{t}$, la traduzione in Java bytecode
semplicemente non riporta esplicitamente il tipo del valore duplicato, dal momento che sar\`a
inferibile. All'interno di \texttt{bytecode/POP.java} scriviamo:
%
\begin{verbatim}
  public InstructionList generateJavaBytecode
      (JavaClassGenerator classGen) {
    // non usiamo il campo type
    return new InstructionList(InstructionFactory.POP);
  }
\end{verbatim}
%
La traduzione del bytecode Kitten $\mathtt{const}\ \mathit{constant}$ \e \piu complessa,
dal momento che ci sono \piu alternative in Java bytecode, che occupano una quantit\`a
variabile di byte. Conseguentemente \e meglio scegliere la versione \piu ottimizzata fra le alternative
disponibili per la specifica $\mathit{constant}$,
come mostrato nella Figura~\ref{fig:kitten_into_java_1}. Dentro
\texttt{bytecode/CONST.java} scriviamo quindi:
%
\begin{verbatim}
  public InstructionList generateJavaBytecode
      (JavaClassGenerator classGen) {
    if (constant == null)
      return new InstructionList
                     (new org.apache.bcel.generic.ACONST_NULL());
    else
      return new InstructionList
                     (classGen.getFactory().createConstant(constant));
  }
\end{verbatim}
%
Con la prima alternativa dell'\texttt{if} gestiamo il caso in cui la costante sia \texttt{nil}.
Con la seconda gestiamo tutti gli altri casi, usando la fattoria BCEL per scegliere il bytecode
\piu ottimizzato per la costante specifica per cui stiamo compilando. Senza l'aiuto di questa fattoria,
avremmo dovuto considerare noi tutti i casi e il relativo bytecode Java \piu ottimizzato per
ciascun caso.

Un altro esempio di utilizzo della fattoria \e per il bytecode Kitten $\mathtt{load}$. Il
corrispondente bytecode Java riporta il numero della variabile locale letta, ma non il suo tipo,
\poiche \e comunque inferibile (Figura~\ref{fig:kitten_into_java_1}). Notiamo qui che
il tipo \texttt{boolean} \e implementato tramite interi (0 significa falso e ogni altro valore vero).
Dentro \texttt{bytecode/LOAD.java} scriviamo:
%
\begin{verbatim}
  public InstructionList generateJavaBytecode
      (JavaClassGenerator classGen) {
    // non usiamo il campo type
    return new InstructionList
                 (InstructionFactory.createLoad(type.toBCEL(), varNum));
  }
\end{verbatim}
%
dove il tipo Kitten viene tradotto nel corrispondente tipo Java bytecode e quindi pasato
alla fattoria per scegliere il bytecode Java corrispondente a quel tipo.

Per adesso a ogni bytecode Kitten \e sempre corrisposta una lista di bytecode Java di lunghezza unitaria.
Nel caso del bytecode Kitten $\mathtt{neg\ boolean}$, la sua traduzione in Java bytecode
\e fatta invece da una lista di 5 istruzioni (Figura~\ref{fig:kitten_into_java_1}). Dentro
\texttt{bytecode/NEG.java} scriviamo:
%
\begin{verbatim}
  public InstructionList generateJavaBytecode
      (JavaClassGenerator classGen) {
    InstructionList il = new InstructionList();

    if (type == BooleanType.INSTANCE) {
      InstructionHandle end = il.insert(InstructionFactory.NOP);
      InstructionHandle after = il.insert(InstructionFactory.ICONST_1);
      il.insert(new org.apache.bcel.generic.GOTO(end));
      il.insert(InstructionFactory.ICONST_0);
      il.insert(new org.apache.bcel.generic.IFEQ(after));
    }
    else
      ((NumericalType) type).neg(il);
    
    return il;
   }
\end{verbatim}
%
Il caso non booleano \e gestito da un metodo \texttt{neg(il)} che aggiunge alla
lista di istruzioni \texttt{il} il bytecode di negazione specifico per quel tipo.
Nel caso della negazione booleana, non esiste un corrispondente bytecode Java di negazione.
Viene costruita invece una lista di 5 bytecode Java che
controllano se il valore in cima allo stack \e falso (\cioe 0) e in tal caso
lo sostituiscono con il valore 1 (\cioe vero); se il valore \e invece vero (diverso da 0)
lo sostituiscono con il valore 0 (\cioe falso). Il metodo BCEL \texttt{insert()}
permette di inserire un bytecode Java all'inizio di una lista di istruzioni.
Tale metodo ritorna un riferimento (\texttt{InstructionHandle}) che \e utile nei casi
in cui si debba costruire un salto a tale punto del codice. Si noti che il riferimento deve
essere ottenuto prima di costruire il salto, per cui il codice viene tipicamente
costruito dall'ultima istruzione verso la prima, come in questo esempio
per $\mathtt{neg\ boolean}$.

Passando alla Figura~\ref{fig:kitten_into_java_2}, \e interessante considerare il caso
della creazione di un oggetto per una costante di tipo stringa. Le stringhe
in Kitten sono implementate dalla classe di supporto
\texttt{runTime/String.java}, per cui
dobbiamo generare del codice che crea un'istanza di tale classe e la inizializza chiamando il
suo costruttore. Dentro \texttt{bytecode/NEWSTRING.java} scriviamo:
%
\begin{verbatim}
  public InstructionList generateJavaBytecode
      (JavaClassGenerator classGen) {

    InstructionFactory factory = classGen.getFactory();
    InstructionList il = new InstructionList();
    String kittenStringName = runTime.String.class.getName();

    il.insert(factory.createInvoke
      (kittenStringName, // class name of the method
      Constants.CONSTRUCTOR_NAME, // name of the method
      org.apache.bcel.generic.Type.VOID, // return type
      new org.apache.bcel.generic.Type[] // parameters types
        { org.apache.bcel.generic.Type.getType("Ljava/lang/String;") },
      Constants.INVOKESPECIAL)); // invokespecial

    il.insert(factory.createConstant(value));
    il.insert(InstructionFactory.DUP);
    il.insert(factory.createNew(kittenStringName));

    return il;
   }
\end{verbatim}
%
Tale codice usa la fattoria BCEL per creare il Java bytecode $\mathtt{new}$ che
crea l'oggetto (in modo da arricchiere la tabella delle costanti),
per creare il bytecode che carica sullo stack il valore della stringa e
per creare il bytecode di invocazione del costruttore. Qui usiamo
il bytecode Java \texttt{invokespecial}, che \e obbligatorio per
invocare i costruttori\footnote{Tale bytecode effettua la ricerca del metodo
a partire dalla classe specificata staticamente nel codice ed \e quindi
ideale anche per la compilazione delle chiamate a un metodo della superclasse,
come con \texttt{super.m(\ldots)} in Java. Nel caso dei costruttori, il
bytecode $\mathtt{invokespecial}$ fallisce se il costruttore non si trova nella classe
indicata, senza effettuare la ricerca nella superclasse, \poiche in Java i
costruttori non si ereditano.}.
